\newpage
\thispagestyle{empty}

\begin{center}
    \Large \textbf{RESUMEN}
\end{center}

    Este proyecto propone el desarrollo de un sistema de soporte a la decisión clínica para oftalmólogos, basado en inteligencia artificial multimodal, que asiste en la detección temprana del queratocono. El sistema procesa de forma automatizada mapas corneales obtenidos con equipos de topografía y tomografía, combinándolos con datos clínicos estructurados del paciente para generar una predicción diagnóstica fundamentada. El prototipo se desarrolla sobre un conjunto de datos público de estudios Orbscan\textsuperscript{\scriptsize\textregistered}{} y se ofrece como una aplicación web bajo un modelo de suscripción.

    El sistema se enmarca como un Software como Producto Médico (SaMD\footnote{Por sus siglas en inglés, \textit{Software as a Medical Device}}) que integra redes neuronales convolucionales para el análisis de imágenes, una rama de procesamiento de datos clínicos estructurados y un mecanismo de fusión multimodal que permite capturar patrones que serían imperceptibles al analizar cada modalidad por separado. Además, incorpora herramientas de explicabilidad (XAI\footnote{Por sus siglas en inglés, \textit{eXplainable Artificial Intelligence}.}) basadas en Grad-CAM que generan mapas de calor sobre las regiones corneales determinantes para cada predicción, proporcionando evidencia visual que fundamenta la decisión y facilita la validación por parte del especialista.

\begin{center}
    \Large \textbf{PALABRAS CLAVE}
\end{center}

    Aplicación web, Explicabilidad (Grad-CAM), Inteligencia artificial multimodal, Queratocono, Redes neuronales convolucionales, Software como Producto Médico (SaMD).

\newpage

